% !PRINTED 482
\weberpage{482}
%<<CONT
\weq{1}{\alpha, \alpha_1, \alpha_2 \ldots \alpha_{n-1}}
sein, die übrigen Wurzeln von $F(x)$
\weq{2}{\alpha_n, \alpha_{n+1} \ldots \alpha_{m-1}.}

Ist $\alpha'$ irgend eine der Grössen (1), $\alpha''$ eine der Grössen (2),
so kann in der Gruppe $P$ keine Permutation vorkommen, durch die $\alpha'$ in
$\alpha''$ übergeführt wird; denn nach Voraussetzung ist
\[
f(\alpha') = 0,
\]
und wenn nun in $P$ eine Permutation vorkäme, durch die $\alpha'$ durch
$\alpha''$ ersetzt würde, so müsste wegen der Eigenschaft \S.\,149, a) auch
$f(\alpha'') = 0$ sein, was unserer Annahme widerspricht.

\gesperrt{Es werden also durch die Permutationen von} $P$ \gesperrt{die
Wurzeln (1) von} $f(x)$ \gesperrt{nur unter einander permutirt.}

Wenn umgekehrt die Gruppe $P$ die Eigenschaft hat, dass ihre Permutationen
einen Theil der $m$ Wurzeln $\alpha$ wie das System (1) nur unter einander
vertauschen, so gestattet das Product
\[
(x - \alpha)(x - \alpha_1) \ldots (x - \alpha_{n-1}) = f(x)
\]
alle Permutationen von $P$, und ist also nach \S.\,149, b) in $\Omega$
enthalten; d. h. $F(x)$ ist reducibel.

Man nennt eine Permutationsgruppe \gesperrt{transitiv}, wenn sie wenigstens
eine Permutation enthält, die ein beliebiges Element in ein beliebiges
anderes überführt; im entgegengesetzten Falle, wenn also die Elemente so in
zwei oder mehr Partien zerlegt werden können, dass durch keine Permutation
der Gruppe ein Element der einen Partie in ein Element der anderen übergeht,
heisst die Gruppe
\gesperrt{intransitiv}. Demnach können wir das Bewiesene in dem Satze
zusammenfassen:

% [sic] Der gedruckte Satz vertauscht die im unmittelbar vorhergehenden Beweis hergeleitete Zuordnung: irreducibel gehört zu transitiv, reducibel zu intransitiv. Der Druckwortlaut bleibt unverändert.
\gesperrt{1. Die Gleichung $F(x) = 0$ ist reducibel oder irreducibel, je
nachdem ihre Galois'sche Gruppe transitiv oder intransitiv ist.}

Wenn ein Theil der Elemente so zusammenhängt, dass durch Permutationen aus
$P$ jedes Element dieses Theiles in jedes andere übergehen kann, so nennen
wir die Elemente dieses Theiles \gesperrt{transitiv verbunden}. Die
verschiedenen transitiv verbundenen Systeme, die durch alle Permutationen von
$P$ nur unter einander permutirt werden, heissen die \gesperrt{Systeme der
Intransitivität}.
